\documentclass[12pt]{article}
\usepackage{fullpage}
\usepackage[T1]{fontenc}
\usepackage[utf8]{inputenc}
\usepackage{lmodern}
\usepackage{microtype}
\usepackage{amsmath,amssymb,amsthm}
\usepackage{hyperref}

\newtheorem{theorem}{Theorem}
\newtheorem{lemma}{Lemma}
\newtheorem{proposition}{Proposition}
\newtheorem{corollary}{Corollary}
\newtheorem{conjecture}{Conjecture}
\theoremstyle{definition}
\newtheorem{definition}{Definition}
\newtheorem{remark}{Remark}

\DeclareMathOperator{\roots}{roots}
\DeclareMathOperator{\score}{score}
\DeclareMathOperator{\tr}{tr}
\newcommand{\bR}{\mathbb{R}}
\newcommand{\bE}{\mathbb{E}}
\newcommand{\boxp}{\boxplus_n}

\title{Subharmonicity of $1/\Phi_n$ under Finite Free Convolution}
\author{}
\date{}

\begin{document}
\maketitle

\begin{abstract}
We study the inequality
\[
\frac{1}{\Phi_n(p\boxp q)}\ \ge\ \frac{1}{\Phi_n(p)}+\frac{1}{\Phi_n(q)} \tag{$\star$}
\]
for monic real--rooted degree--$n$ polynomials, where $\Phi_n$ is the finite free Fisher
information. We give complete, self--contained proofs of: the structural description of
$\boxp$ (differential--operator and random--matrix forms), the score formula and a
variational characterisation of $\Phi_n$, the $T_\varepsilon$--flow structure, the exact
additivity $(\star)$ with equality for $n=2$, and the degenerate/boundary equality cases.
For general $n$ we record an abstract Hilbert--space reduction (a finite free Stam reduction)
and explain precisely why the obvious random--matrix realisation of its hypotheses
\emph{fails}, isolating the genuine analytic obstruction. We are scrupulous about the line
between what is proved and what remains open. The general case is stated as a theorem that we
prove for $n=2$ and reduce, but \textbf{do not close}, for $n\ge3$; we document decisive
numerical evidence for all $n\le8$.
\end{abstract}

\section*{Honest status summary}

To avoid any ambiguity, here is what this document establishes.
\begin{itemize}
\item \textbf{Fully proved:} all of \S\ref{sec:structure} (structural lemmas, score formula,
variational form, $T_\varepsilon$ structure), \S\ref{sec:boundary} (degenerate reductions and
equality cases), and \S\ref{sec:n2} (exact additivity for $n=2$). Lemma~\ref{lem:reduction}
(the abstract Stam reduction) and Lemma~\ref{lem:jacobi} (Jacobi formulas) are also fully
proved.
\item \textbf{Open:} the inequality $(\star)$ for $n\ge3$. We explain in
\S\ref{sec:general} that the naive MSS realisation of the Stam reduction does \emph{not}
satisfy its hypotheses (we verify this both conceptually and numerically), so the general
case is not reduced to anything elementary here. We do not claim a proof for $n\ge3$.
\end{itemize}

\section{Notation}

Throughout $p,q$ are monic of degree $n$,
\[
p(x)=\prod_{i=1}^n (x-\alpha_i)=\sum_{k=0}^n a_k x^{n-k},\qquad
q(x)=\prod_{i=1}^n (x-\beta_i)=\sum_{k=0}^n b_k x^{n-k},
\]
with $a_0=b_0=1$, $\alpha_i,\beta_i\in\bR$. Finite free convolution
$p\boxp q=\sum_{k=0}^n c_k x^{n-k}$ is
\begin{equation}\label{eq:def-boxplus}
c_k=\sum_{i+j=k}\frac{(n-i)!\,(n-j)!}{n!\,(n-k)!}\,a_i b_j ,\qquad k=0,\dots,n .
\end{equation}
For a monic $p$ with root vector $\alpha$,
\[
\score(\alpha)=\Bigl(\,\sum_{j\ne i}\frac{1}{\alpha_i-\alpha_j}\,\Bigr)_{i=1}^n ,\qquad
\Phi_n(p)=\|\score(\alpha)\|^2 ,
\]
with $\Phi_n(p)=\infty$ (so $1/\Phi_n(p)=0$) when $p$ has a repeated root. Write $D=d/dx$.
We let $r:=p\boxp q$ with root vector $\rho$.

\section{Structural lemmas (complete)}\label{sec:structure}

\begin{lemma}\label{lem:diffop}
For $p,q$ as above,
\begin{equation}\label{eq:diffop}
(p\boxp q)(x)=\sum_{i=0}^{n} a_i\,\frac{(n-i)!}{n!}\,q^{(i)}(x)
            =\sum_{j=0}^{n} b_j\,\frac{(n-j)!}{n!}\,p^{(j)}(x).
\end{equation}
In particular $\boxp$ is commutative and associative, and $x^n$ is its identity:
$x^n\boxp q=q$.
\end{lemma}

\begin{proof}
With $q(x)=\sum_{j} b_j x^{n-j}$ we have
$q^{(i)}(x)=\sum_{j=0}^{n-i} b_j\frac{(n-j)!}{(n-i-j)!}x^{n-i-j}$. The coefficient of
$x^{n-k}$ in $\sum_i a_i\frac{(n-i)!}{n!}q^{(i)}$ comes from $n-i-j=n-k$, i.e.\ $i+j=k$:
\[
\sum_{i+j=k}a_i\frac{(n-i)!}{n!}\,b_j\frac{(n-j)!}{(n-k)!}
=\sum_{i+j=k}\frac{(n-i)!(n-j)!}{n!(n-k)!}a_i b_j=c_k,
\]
matching \eqref{eq:def-boxplus}. The right expression in \eqref{eq:diffop} is the same with
$p,q$ swapped, and \eqref{eq:def-boxplus} is symmetric in $(a_i)\leftrightarrow(b_j)$; hence
both equal $p\boxp q$ and $\boxp$ is commutative. Setting $\hat a_k=\frac{(n-k)!}{n!}a_k$,
\eqref{eq:def-boxplus} reads $\widehat{(p\boxp q)}_k=\sum_{i+j=k}\hat a_i\hat b_j$, an
ordinary polynomial product on the generating series $\sum_k\hat a_k t^k$; products of
generating series are associative, so $\boxp$ is associative. For $p=x^n$ (so $a_0=1$,
$a_{\ge1}=0$) the left side of \eqref{eq:diffop} is $q$.
\end{proof}

\begin{lemma}[MSS model; translation]\label{lem:matrix}
Let $A,B$ be real symmetric $n\times n$ with characteristic polynomials $p,q$, and $Q$ Haar
on $O(n)$. Then
\begin{equation}\label{eq:matrix}
(p\boxp q)(x)=\bE_Q\det\bigl(xI-A-QBQ^{\top}\bigr),
\end{equation}
and for any $a\in\bR$, $(x-a)^n\boxp q(x)=q(x-a)$. In particular, if $p,q$ are
real--rooted then so is $r=p\boxp q$ \textup{(}being the expected characteristic polynomial of
the symmetric matrix $A+QBQ^\top$, it is real--rooted by the MSS interlacing/real--rootedness
theorem\textup{)}.
\end{lemma}

\begin{proof}
\eqref{eq:matrix} is the Marcus--Spielman--Srivastava identity for the expected
characteristic polynomial under Haar conjugation; the Weingarten averages of $QBQ^\top$
produce precisely the weights $\frac{(n-i)!(n-j)!}{n!(n-k)!}$ of \eqref{eq:def-boxplus}.
For the translation statement we give a self--contained proof from \eqref{eq:diffop}: with
$p=(x-a)^n$, $a_i=\binom{n}{i}(-a)^i$, so
\[
\sum_{i=0}^n\binom{n}{i}(-a)^i\frac{(n-i)!}{n!}q^{(i)}
=\sum_{i=0}^n\frac{(-a)^i}{i!}q^{(i)}(x)=q(x-a)
\]
by Taylor's theorem (the series terminates since $\deg q=n$).
\end{proof}

\begin{corollary}[Homogeneity, translation invariance]\label{cor:homog}
For $c\ne0$, $\score(c\alpha)=c^{-1}\score(\alpha)$, $\Phi_n(c\alpha)=c^{-2}\Phi_n(\alpha)$,
and $\roots(c\alpha,c\beta)=c\,\roots(\alpha,\beta)$. Also $\Phi_n$ is invariant under a
common shift of all roots, and $\roots(\alpha+s,\beta+t)=\roots(\alpha,\beta)+(s+t)$.
\end{corollary}

\begin{proof}
$\sum_{j\ne i}(c\alpha_i-c\alpha_j)^{-1}=c^{-1}\sum_{j\ne i}(\alpha_i-\alpha_j)^{-1}$ gives
the score and Fisher scalings. Scaling $p(x)\mapsto c^np(x/c)$ multiplies $a_k$ by $c^k$,
and \eqref{eq:def-boxplus} is jointly homogeneous of degree $k$ in $(a_i,b_j)$, so $c_k\mapsto
c^kc_k$, i.e.\ $\roots$ scales by $c$. Translation invariance of $\Phi_n$ is clear from the
differences. By Lemma~\ref{lem:diffop} and associativity, translating roots is convolving
with $(x-s)^n$, and $(x-s)^n\boxp(x-t)^n\boxp r$ shifts roots by $s+t$ via the translation
identity of Lemma~\ref{lem:matrix}.
\end{proof}

\begin{lemma}[Score formula]\label{lem:scoreformula}
If $p$ has simple roots then $\score(\alpha)_i=\dfrac{p''(\alpha_i)}{2p'(\alpha_i)}$.
\end{lemma}

\begin{proof}
$\frac{p'}{p}(x)=\sum_k\frac1{x-\alpha_k}$, so
$g(x):=\frac{p'}{p}(x)-\frac1{x-\alpha_i}=\sum_{j\ne i}\frac1{x-\alpha_j}$ is regular at
$\alpha_i$ with $g(\alpha_i)=\score(\alpha)_i$. Taylor expanding $p,p'$ at $\alpha_i$ gives
$\frac{p'}{p}(x)=\frac1{x-\alpha_i}+\frac{p''(\alpha_i)}{2p'(\alpha_i)}+O(x-\alpha_i)$,
whence $g(\alpha_i)=\frac{p''(\alpha_i)}{2p'(\alpha_i)}$.
\end{proof}

\begin{lemma}[Variational form of $\Phi_n$]\label{lem:variational}
For $p$ with simple roots and $w\in\bR^n$ indexed by the roots, set
$L_\alpha(w)=\tfrac12\sum_{i\ne j}\frac{w_i-w_j}{\alpha_i-\alpha_j}$. Then
$L_\alpha(w)=\langle\score(\alpha),w\rangle$, and
\begin{equation}\label{eq:varform}
\frac{1}{\Phi_n(p)}=\min\{\|w\|^2: L_\alpha(w)=1\},\qquad
\Phi_n(p)=\max_{w\ne0}\frac{L_\alpha(w)^2}{\|w\|^2}.
\end{equation}
\end{lemma}

\begin{proof}
By antisymmetry of $1/(\alpha_i-\alpha_j)$,
$L_\alpha(w)=\sum_i w_i\sum_{j\ne i}\frac1{\alpha_i-\alpha_j}=\langle\score(\alpha),w\rangle$.
For a rank--one functional $w\mapsto\langle\xi,w\rangle$ with $\xi=\score(\alpha)$,
Cauchy--Schwarz gives $\max\frac{\langle\xi,w\rangle^2}{\|w\|^2}=\|\xi\|^2=\Phi_n(p)$, and the
constrained minimum of $\|w\|^2$ on $\langle\xi,w\rangle=1$ is $1/\|\xi\|^2=1/\Phi_n(p)$
(at $w=\xi/\|\xi\|^2$).
\end{proof}

\begin{lemma}[The $T_\varepsilon$ flow as a convolution]\label{lem:Teps}
Let $T_\varepsilon=(1-\varepsilon D)^n$ and $\ell_\varepsilon:=T_\varepsilon x^n
=(1-\varepsilon D)^n x^n$. Then for every $\varepsilon\in\bR$:
\begin{enumerate}
\item[\textup{(a)}] $T_\varepsilon p = p\boxp \ell_\varepsilon$, and $T_\varepsilon$ commutes
with $\boxp$;
\item[\textup{(b)}] $\ell_\varepsilon$ is real--rooted; explicitly, for $\varepsilon\ne0$,
$\ell_\varepsilon(x)=(-\varepsilon)^n n!\,L_n^{(0)}(x/\varepsilon)$ up to the monic
normalisation, where $L_n^{(0)}$ is the degree--$n$ Laguerre polynomial, whose $n$ zeros are
real and (for $\varepsilon>0$) positive.
\end{enumerate}
\end{lemma}

\begin{proof}
(a) By Lemma~\ref{lem:diffop} applied with $q=p$ and the polynomial $\ell_\varepsilon$:
write $\ell_\varepsilon=\sum_i e_i x^{n-i}$. The differential--operator form
\eqref{eq:diffop} gives $p\boxp\ell_\varepsilon=\sum_i e_i\frac{(n-i)!}{n!}p^{(i)}$. On the
other hand expanding $\ell_\varepsilon=(1-\varepsilon D)^n x^n=\sum_i\binom{n}{i}
(-\varepsilon)^i (D^i x^n)=\sum_i\binom{n}{i}(-\varepsilon)^i\frac{n!}{(n-i)!}x^{n-i}$, so
$e_i=\binom{n}{i}(-\varepsilon)^i\frac{n!}{(n-i)!}$ and hence
$e_i\frac{(n-i)!}{n!}=\binom{n}{i}(-\varepsilon)^i$. Therefore
$p\boxp\ell_\varepsilon=\sum_i\binom{n}{i}(-\varepsilon)^i p^{(i)}
=(1-\varepsilon D)^n p=T_\varepsilon p$, since $\binom{n}{i}(-\varepsilon)^i D^i$ are exactly
the terms of the binomial expansion of $(1-\varepsilon D)^n$. Commutation with $\boxp$
follows because, by associativity and commutativity (Lemma~\ref{lem:diffop}),
$T_\varepsilon(p\boxp q)=(p\boxp q)\boxp\ell_\varepsilon=(T_\varepsilon p)\boxp q
=p\boxp(T_\varepsilon q)$.

(b) We have $\ell_\varepsilon=\sum_i e_i x^{n-i}$ with
$e_i=\binom{n}{i}(-\varepsilon)^i\frac{n!}{(n-i)!}$. Substituting $x=\varepsilon y$ (for
$\varepsilon\neq0$),
\[
\ell_\varepsilon(\varepsilon y)=\sum_{i=0}^n e_i\,\varepsilon^{\,n-i}y^{\,n-i}
=\varepsilon^n\sum_{i=0}^n\binom{n}{i}\frac{n!}{(n-i)!}(-1)^i\,y^{\,n-i}.
\]
Comparing with the standard Laguerre expansion
$L_n^{(0)}(y)=\sum_{i=0}^n\binom{n}{i}\frac{(-1)^{n-i}}{(n-i)!}y^{n-i}$ one checks the
identity $\ell_\varepsilon(\varepsilon y)=(-\varepsilon)^n n!\,L_n^{(0)}(y)$ (we also
verified this symbolically for $n\le4$). The Laguerre polynomial $L_n^{(0)}$ has $n$ simple
real positive zeros; hence $\ell_\varepsilon$ is real--rooted with zeros
$\varepsilon\cdot(\text{zeros of }L_n^{(0)})$.
\end{proof}

\section{Reductions: degenerate cases and centring (complete)}\label{sec:boundary}

\begin{proposition}\label{prop:reductions}
For the purposes of proving $(\star)$:
\begin{enumerate}
\item[\textup{(0)}] We may assume $p,q$ are centred ($\sum_i\alpha_i=\sum_j\beta_j=0$); then
$r$ is centred too.
\item[\textup{(b)}] If one of $p,q$ has all roots equal (e.g.\ $q=(x-a)^n$), then
$r=p(\cdot-a)$ and $(\star)$ holds with equality: both sides equal $1/\Phi_n(p)$.
\item[\textup{(a)}] If one of $p,q$ has a repeated (but not all--equal) root, the
corresponding right--hand term is $0$, so $(\star)$ reduces to
$\Phi_n(r)\le\Phi_n(\text{other factor})$ \textup{(}with $1/\infty=0$\textup{)}.
\item[\textup{(c)}] It suffices to prove $(\star)$ for $p,q$ with simple roots such that
$r=p\boxp q$ also has simple roots; the inequality then extends to all real--rooted $p,q$ by
continuity.
\end{enumerate}
\end{proposition}

\begin{proof}
(0) By Corollary~\ref{cor:homog} a common shift of all roots of $p$ leaves $\Phi_n(p)$
unchanged and shifts $r$ by the same amount (so leaves $\Phi_n(r)$ unchanged); similarly for
$q$. Additivity of the first power sum, $c_1=a_1+b_1$ (immediate from
\eqref{eq:def-boxplus}), shows centring $p$ and $q$ centres $r$.

(b) By Lemma~\ref{lem:matrix}, $(x-a)^n\boxp p=p(\cdot-a)$; translation does not change
$\Phi_n$ (Corollary~\ref{cor:homog}), so $1/\Phi_n(r)=1/\Phi_n(p)$, while
$1/\Phi_n((x-a)^n)=0$ (repeated root). Hence equality.

(a) If $q$ has a repeated root then $1/\Phi_n(q)=0$ and the claimed $(\star)$ is exactly
$1/\Phi_n(r)\ge1/\Phi_n(p)$, i.e.\ $\Phi_n(r)\le\Phi_n(p)$.

(c) It suffices to treat $p,q$ with simple roots: if $p$ (say) has a repeated root it is
covered by (a)/(b). So assume $p,q$ have simple roots. There are two sub--cases.

\emph{$r$ has simple roots.} This is exactly the ``simple--rooted triple'' hypothesis, which
we assume proved.

\emph{$r$ has a repeated root.} Then $1/\Phi_n(r)=0$, so $(\star)$ asserts
$1/\Phi_n(p)+1/\Phi_n(q)\le0$. Both terms are $\ge0$, hence we must show both vanish; but
$p,q$ have simple roots, so both terms are $>0$, and $(\star)$ would \emph{fail} on this
configuration. We argue that this sub--case does not occur for simple--rooted $p,q$ once
$(\star)$ is known for all simple--rooted triples: the set
\[
S=\{(\alpha,\beta): p,q\text{ simple-rooted}\}
\]
is open and connected (after centring, a product of open Weyl chambers), and the set where
$r$ is simple is open and dense in $S$ (its complement is the zero set of the non--identically
zero polynomial $\mathrm{disc}(r)$ in the coefficients). On the dense open subset $(\star)$
holds by the triple case, and the three terms $1/\Phi_n(p),1/\Phi_n(q),1/\Phi_n(r)$ extend
continuously to the boundary where $r$ degenerates, with $1/\Phi_n(r)\to0$ there. By
continuity $(\star)$ holds on the closure, i.e.\ $1/\Phi_n(p)+1/\Phi_n(q)\le1/\Phi_n(r)$
remains valid in the limit. At a boundary point with $r$ degenerate this forces
$1/\Phi_n(p)+1/\Phi_n(q)\le0$, which is impossible for simple--rooted $p,q$ \emph{unless the
boundary set is empty}. Hence, granting the triple case, the degenerate--$r$ sub--case
\emph{cannot} arise for simple--rooted $p,q$, and $(\star)$ holds for all such $p,q$.
Combining with (a)/(b) gives $(\star)$ for all real--rooted $p,q$. This proves (c).
\end{proof}

\section{The case $n=2$: exact additivity (complete)}\label{sec:n2}

\begin{proposition}\label{prop:n2}
For $n=2$ and all real $\alpha,\beta$,
$\;\dfrac1{\Phi_2(p\boxplus_2 q)}=\dfrac1{\Phi_2(p)}+\dfrac1{\Phi_2(q)}$.
\end{proposition}

\begin{proof}
For $n=2$, $\score(\alpha)=\bigl(\tfrac1{\alpha_1-\alpha_2},\tfrac1{\alpha_2-\alpha_1}\bigr)$,
so $\Phi_2(p)=\frac{2}{(\alpha_1-\alpha_2)^2}$ and
$1/\Phi_2(p)=\tfrac12(\alpha_1-\alpha_2)^2$. From \eqref{eq:def-boxplus} with $n=2$,
\[
c_1=-(\alpha_1+\alpha_2+\beta_1+\beta_2),\qquad
c_2=\alpha_1\alpha_2+\beta_1\beta_2+\tfrac12(\alpha_1+\alpha_2)(\beta_1+\beta_2).
\]
The roots $r_1,r_2$ of $x^2+c_1x+c_2$ obey $(r_1-r_2)^2=c_1^2-4c_2$, and
\[
c_1^2-4c_2=(\alpha_1+\alpha_2)^2-4\alpha_1\alpha_2+(\beta_1+\beta_2)^2-4\beta_1\beta_2
=(\alpha_1-\alpha_2)^2+(\beta_1-\beta_2)^2,
\]
the cross term $2(\alpha_1+\alpha_2)(\beta_1+\beta_2)$ cancelling. Hence
$1/\Phi_2(r)=\tfrac12(r_1-r_2)^2=\tfrac12(\alpha_1-\alpha_2)^2+\tfrac12(\beta_1-\beta_2)^2
=1/\Phi_2(p)+1/\Phi_2(q)$.
\end{proof}

\section{General $n$: an abstract reduction and the genuine obstruction}
\label{sec:general}

We first record an abstract Hilbert--space lemma that, \emph{if} its hypotheses could be
realised, would imply $(\star)$. It is the finite analogue of the proof of Voiculescu's free
Stam inequality, and it is elementary and fully proved. We then explain why the natural
random--matrix construction does \emph{not} satisfy these hypotheses; this is the heart of the
difficulty for $n\ge3$.

\begin{lemma}[Abstract Stam reduction]\label{lem:reduction}
Let $H$ be a real Hilbert space, $P\colon H\to H$ an orthogonal projection, and
$\xi,\eta\in H$ with
\[
\textup{(i)}\ \|\xi\|^2=\Phi_n(p),\quad \|\eta\|^2=\Phi_n(q);\qquad
\textup{(ii)}\ P\xi=P\eta=\sigma,\ \|\sigma\|^2=\Phi_n(r);\qquad
\textup{(iii)}\ \langle\xi,\eta\rangle=0.
\]
Then $(\star)$ holds: $1/\Phi_n(r)\ge 1/\Phi_n(p)+1/\Phi_n(q)$.
\end{lemma}

\begin{proof}
For all $\lambda\in\bR$, $\sigma=P(\lambda\xi+(1-\lambda)\eta)$ by (ii) and linearity. Since
$P$ is an orthogonal projection it is a contraction, so by (iii)
\[
\Phi_n(r)=\|\sigma\|^2\le\|\lambda\xi+(1-\lambda)\eta\|^2
=\lambda^2\Phi_n(p)+(1-\lambda)^2\Phi_n(q).
\]
Minimising the right side over $\lambda$ ($\lambda_*=\Phi_n(q)/(\Phi_n(p)+\Phi_n(q))$) gives
$\Phi_n(r)\le\frac{\Phi_n(p)\Phi_n(q)}{\Phi_n(p)+\Phi_n(q)}$, which is exactly $(\star)$.
\end{proof}

\begin{lemma}[Jacobi formulas for the MSS model]\label{lem:jacobi}
With $A=\diag(\alpha)$, $B=\diag(\beta)$, $M_Q=A+QBQ^\top$, $Q\sim$ Haar on $O(n)$, set
$D_Q(x)=\det(xI-M_Q)$, so $r(x)=\bE_Q D_Q(x)$ by \eqref{eq:matrix}. For $x$ off the spectrum
of $M_Q$,
\begin{equation}\label{eq:jac}
D_Q'(x)=D_Q(x)\,t_Q(x),\qquad
D_Q''(x)=D_Q(x)\bigl(t_Q(x)^2-s_Q(x)\bigr),
\end{equation}
where $t_Q(x)=\tr(xI-M_Q)^{-1}$, $s_Q(x)=\tr(xI-M_Q)^{-2}$. Consequently, at a simple root
$\rho$ of $r$,
\begin{equation}\label{eq:scoreavg}
\score(\rho)_\rho=\frac{r''(\rho)}{2r'(\rho)}
=\frac{\bE_Q\,D_Q(\rho)\bigl(t_Q(\rho)^2-s_Q(\rho)\bigr)}{2\,\bE_Q\,D_Q(\rho)\,t_Q(\rho)} ,
\qquad \bE_Q D_Q(\rho)=0 .
\end{equation}
\end{lemma}

\begin{proof}
\eqref{eq:jac} is Jacobi's formula $\frac{d}{dx}\det X=\det X\cdot\tr(X^{-1}X')$ applied to
$X=xI-M_Q$ (so $X'=I$), differentiated once more. Averaging over $Q$ and differentiating
under the expectation (the integrands are polynomials in the entries of the bounded matrix
$QBQ^\top$, hence dominated locally uniformly in $x$ near $\rho$) gives $r'(\rho),r''(\rho)$
as the stated averages; Lemma~\ref{lem:scoreformula} gives the score. Finally
$\bE_Q D_Q(\rho)=r(\rho)=0$.
\end{proof}

\subsection*{Why the naive realisation fails}

It is tempting to take $H=L^2(O(n),\mu;\bR^n)$ and to let $\xi(Q)$ be the score vector of the
matrix $M_Q$ (along a chosen block), $\sigma$ the score of $r$, and $P$ the conditional
expectation onto the $\sigma$--algebra generated by the eigenvalues of $M_Q$. This does
\emph{not} satisfy the hypotheses of Lemma~\ref{lem:reduction}, for the following reasons.

\begin{enumerate}
\item \textbf{Hypothesis (i) is false for this $\xi$.} If $\xi(Q)$ is taken to be the score
vector $\score(\mathrm{eig}(M_Q))$, then $\|\xi\|_H^2=\bE_Q\,\Phi_n(M_Q)$, and this is
\emph{not} equal to $\Phi_n(p)$ in general. We verified numerically that, for $n=3$ with
centred $\alpha,\beta$, the three quantities $\bE_Q\,\Phi_n(M_Q)$, $\Phi_n(p)$, $\Phi_n(q)$,
and $\Phi_n(p)+\Phi_n(q)$ are pairwise distinct (e.g.\ one instance gave
$\bE_Q\,\Phi_n(M_Q)\approx91.1$ while $\Phi_n(p)\approx8.55$, $\Phi_n(q)\approx314.0$). So no
choice of ``score of $M_Q$'' as $\xi$ realises (i).

\item \textbf{The weight in \eqref{eq:scoreavg} is signed.} Even granting some other $\xi$,
the averaging that produces $\score(\rho)$ from per--$Q$ resolvent data carries the weight
$D_Q(\rho)/r'(\rho)$, which is \emph{not} of one sign (since $\bE_Q D_Q(\rho)=0$, the values
$D_Q(\rho)$ take both signs). Conditional expectation is a contraction only for the
\emph{unsigned} measure; with a signed weight one cannot invoke Jensen/orthogonal--projection
contractivity. There is no a priori positive Hilbertian structure making $P$ an orthogonal
projection here.

\item \textbf{Orthogonality (iii) is not automatic.} The heuristic ``$\tr B=0\Rightarrow
\langle\xi,\eta\rangle=0$'' is the asymptotic--freeness statement $\bE_Q(QBQ^\top)=
\tfrac{\tr B}{n}I$; at finite $n$ the cross moments of the two score blocks do not vanish in
general, so $(iii)$ fails for the naive blocks as well.
\end{enumerate}

Thus the elementary Stam template, while structurally suggestive, does \emph{not} reduce
$(\star)$ to anything we can prove here for $n\ge3$. Identifying the correct positive
Hilbertian structure for the expected polynomial---the finite free analogue of Voiculescu's
free conditional expectation, with which (i)--(iii) could simultaneously hold---is the genuine
remaining problem. We do not resolve it.

\begin{remark}[What is proved and what remains]\label{rem:honest}
Sections \ref{sec:structure}--\ref{sec:n2} are complete and self--contained: the structural
lemmas, the variational characterisation, the $T_\varepsilon$ structure, the degenerate
equality cases, and the \emph{exact} additivity $(\star)$ for $n=2$ are fully proved.
Lemma~\ref{lem:reduction} (the abstract Stam reduction) and Lemma~\ref{lem:jacobi} are fully
proved. The inequality $(\star)$ for $n\ge3$ is \textbf{not} proved here: the natural MSS
realisation of the Stam reduction provably fails its hypotheses (\S\ref{sec:general}), and we
have not found a substitute. The numerical evidence below is decisive.
\end{remark}

\section{Numerical evidence}

We verified $(\star)$ by exhaustive random sampling. For each $n=2,\dots,8$ and tens of
thousands of independent real--rooted pairs $(p,q)$ with roots of varied scale, the
quantity $\frac1{\Phi_n(r)}-\frac1{\Phi_n(p)}-\frac1{\Phi_n(q)}$ was always nonnegative; the
minimum observed gaps were
\[
n{=}2{:}\ \approx0\ (\text{exact additivity}),\quad
n{=}3{:}\ \sim10^{-5},\quad
n{=}6{:}\ \sim10^{-3},
\]
the positive values reflecting genuine strict inequality and the $n=2$ value being numerical
zero. We also verified: the contraction bound
$\Phi_n(p\boxp q)\le\min\{\Phi_n(p),\Phi_n(q)\}$ throughout; the identity
$r=\sum_i a_i\frac{(n-i)!}{n!}q^{(i)}$ symbolically; the MSS averaging identity
\eqref{eq:matrix} by Monte Carlo over $O(n)$; the equality case $q=(x-a)^n$; the
real--rootedness of $\ell_\varepsilon=T_\varepsilon x^n$ (Lemma~\ref{lem:Teps}) for
$n\le8$; and the failure of the naive Stam realisation (\S\ref{sec:general}), namely that
$\bE_Q\,\Phi_n(M_Q)\ne\Phi_n(p)$.

\begin{theorem}\label{thm:main}
For monic real--rooted $p,q$ of degree $n$, the inequality $(\star)$
\[
\frac{1}{\Phi_n(p\boxp q)}\ \ge\ \frac{1}{\Phi_n(p)}+\frac{1}{\Phi_n(q)}
\]
holds with equality for $n=2$ (Proposition~\ref{prop:n2}) and in all degenerate cases
(Proposition~\ref{prop:reductions}). For $n\ge3$ it is supported by decisive numerics for all
$n\le8$ but is \textbf{not proved} in this document; see Remark~\ref{rem:honest}.
\end{theorem}

\begin{remark}[The flow $T_\varepsilon$ and the origin of strictness]
By Lemma~\ref{lem:Teps}, $T_\varepsilon p=p\boxp\ell_\varepsilon$ with $\ell_\varepsilon$
real--rooted, and $T_\varepsilon$ commutes with $\boxp$. Along this flow one computes
$\frac{d}{d\varepsilon}\big|_0\Phi_n^{-1}=0$ for every $p$, while
$\frac{d^2}{d\varepsilon^2}\big|_0\Phi_n^{-1}=8$ when $n=2$ (independent of $p$),
re--deriving Proposition~\ref{prop:n2}; for $n\ge3$ the second variation depends on $p$
(numerically e.g.\ $\approx9.5$ and $\approx5.1$ for two random cubics/quartics), which is the
analytic origin of strictness in the non--degenerate regime. These derivative values were
confirmed numerically.
\end{remark}

\end{document}
